% answers_short.tex —— 1977 年·江苏无锡市初试卷 · 简短答案速查
% 由 tools/make_short_answers.py 自 answers.tex 的 \gksolhead/\gkanswer 提取，
% 并与 1977/1977_数学真题_江苏无锡市初试_答案要点.md 逐题比对后生成，勿手改（改源 MD 后重跑）。

\begin{gknotice}
\textbf{使用说明}：本节只给最终答案便于快速核对，完整推导见解析版。答案为编写者独立推导，\textbf{非官方标准答案}。
\end{gknotice}

\gkpart{题目}
\begin{itemize}[leftmargin=2.6em]
  \item[\textbf{1.}] (1) $f(0)=-5$；$f(-1)=0$；$f\left(\dfrac13\right)=-\dfrac{16}{3}$；(2) $\sin A=\dfrac5{13}$；$\cos B=\dfrac5{13}$；$\operatorname{tg}A=\dfrac5{12}$；$\operatorname{ctg}A=\dfrac{12}{5}$；(3) $\sqrt{(3-x)^2}=\left|3-x\right|$（$x\le 3$ 时为 $3-x$，$x>3$ 时为 $x-3$）；(4) 极坐标方程 $\rho=5$；圆心的极坐标为 $(0,0)$（圆心即极点，极角任意）；(5) 体积 $30\ \mathrm{cm^3}$；表面积 $62\ \mathrm{cm^2}$
  \item[\textbf{2.}] (1) $-1$；(2) $\dfrac{a}{b^2}$；(3) $4$；(4) $\dfrac{\sqrt6}{2}$；(5) $\dfrac34$；(6) $\dfrac{2}{\left|\sin\theta\right|}$
  \item[\textbf{3.}] (1) $x=-\dfrac25$；(2) $x=4$，$y=3$；(3) 有两个相等的实数根时 $m=\dfrac43$；无实数根时 $0<m<\dfrac43$
  \item[\textbf{4.}] 两数列通项分别为 $a_n=2n-1$ 与 $b_n=\dfrac{1}{2^{n-1}}$；前 $10$ 项之和 $S_{10}=\dfrac{52223}{512}\approx 102$。
  \item[\textbf{5.}] $a=\sqrt6$，$B=75^\circ$，$C=45^\circ$。
  \item[\textbf{6.}] $R=a$。
  \item[\textbf{7.}] 直线方程为 $3x-4y+12=0$（过焦点 $(0,3)$）或 $3x-4y-12=0$（过焦点 $(0,-3)$）；直线与坐标轴围成的三角形面积为 $6$。
  \item[\textbf{8.}] 见解析版
  \item[\textbf{9.}] 见解析版
\end{itemize}
